\begin{figure}[t]
\centering
\resizebox{\textwidth}{!}{%
\begin{tikzpicture}[
  >=Latex,
  node distance=1.25cm,
  box/.style={draw, rounded corners=2pt, align=center, minimum height=1.55cm, minimum width=3.1cm, inner sep=6pt},
  arr/.style={->, line width=0.8pt},
  lab/.style={align=center, font=\small}
]

\node[box] (source) {\textbf{Equal-charge local code}\\
$V_{\mathcal C}^\dagger Q_A V_{\mathcal C}=q_A I$\\
conserved source; handoff $T_*$};

\node[box, right=of source] (gmode) {\textbf{Normalized graviton mode}\\
complete encoder + tail waveform\\
$\int |f(t)|^2dt=1$};

\node[box, right=of gmode] (rport) {\textbf{Receiver gravitational port}\\
retarded wave-zone access\\
$R/c$ causal lower bound};

\node[box, right=of rport] (memory) {\textbf{Receiver memory}\\
$\tau_c=\eta_Q^{\rm link}\mathcal T_f$\\
$\Delta_{\rm mem}=\tau_c-m_c$};

\node[box, right=of memory] (readout) {\textbf{Accessible register}\\
readout $(\tau_r,m_r)$\\
$\Delta_{\rm acc}=\tau_r\Delta_{\rm mem}-m_r$};

\draw[arr] (source) -- node[lab, above] {$\displaystyle \beta_{g,A}=\frac{\kappa_{g,A}}{\kappa_A}$\\source branching} (gmode);

\draw[arr] (gmode) -- node[lab, above] {$\displaystyle \eta_{\rm store}=\frac{25\mathcal O}{16(kR)^2}$\\mode capture} (rport);

\draw[arr] (rport) -- node[lab, above] {$\displaystyle \beta_{g,B}\mathcal T_f$\\branching $\times$ loading} (memory);

\draw[arr] (memory) -- node[lab, above] {noisy readout} (readout);

\node[lab, below=0.55cm of gmode] (backbone) {$\displaystyle \eta_Q^{\rm link}=\beta_{g,A}\eta_{\rm store}\beta_{g,B}$};

\end{tikzpicture}%
}
\caption{Source-resolved gravitational quantum-link architecture after the local code is instantiated. The three multiplicative factors $\beta_{g,A}$, $\eta_{\rm store}$, and $\beta_{g,B}$ form the coherent link backbone; the fixed waveform contributes $\mathcal T_f\le1$. Noise determines the memory quantum excess, and a separate readout stage determines whether that excess remains accessible. Encoder precursor radiation is retained in the complete waveform $f(t)$ rather than reset at the handoff.}
\label{fig:link-budget}
\end{figure}
